% Local input for fourfolds-owned D.2.45. No separate theorem ID.
\subsubsection*{The original type-II fibre and its good logarithmic replacement}

This input concerns the original type-$II$ cap and the actual free
good-reduction logarithmic cap of multiplicity $m=6k$.  It does not
assert the same formula for star fibres, nonfree quotients, or a
nonreduced model without first removing its divisorial differential
factor.  All statements are about the smooth total threefold.
In the profiles in question the linearization has its zero at
a different, multiplicative fibre; it is a unit at the type-$II$
site.  No collision of that zero with the additive cusp is
being included.

\paragraph{The original cap: coefficient two.}
At the cusp of a type-$II$ elliptic fibre the total elliptic surface
is smooth and its projection has an ordinary $A_2$ plane-curve
critical point.  Locally its equation can be put in the form
\[
                         q=y^2-x^3.
 \tag{D.IIF1}
\]
For example the Weierstrass equation with $b'(0)\ne0$ solves for
the base coordinate by the implicit-function theorem.  The
holomorphic splitting lemma in the nondegenerate $y$ direction
then writes the map as $Y^2+h(X)$, with $h$ of order three;
taking the holomorphic cube root of its unit coefficient gives
(D.IIF1).  Thus this assertion is a statement about the map to
the base, not only about the singular zero set.

The line-coordinate construction of the original cap is locally
a product with a nonvanishing scalar coordinate.  Its covering
charts are therefore smooth coordinates $(x,y,z)$ with (D.IIF1)
independent of $z$.  Their identifications are holomorphic local
isomorphisms preserving the base map.  At the singular point of
the elliptic fibre, the two scalar period directions form the
compact critical elliptic curve $E_{II}$.
The differential and its kernel have the explicit form
\[
 df=(-3x^2,2y,0),\qquad
 V=\mathcal O(2y,3x^2,0)\oplus\mathcal O(0,0,1).
 \tag{D.IIF2}
\]
The two coefficients $x^2,y$ are relatively prime, so these
generators give the whole kernel.  In particular $V$ is locally
free, and the scalar vector $\delta$ is a nonzero multiple of
the second abstract generator, including on $E_{II}$.
The cokernel of $df$ is $L\otimes\mathcal O/(x^2,y)$.
At the generic point of $E_{II}$ this scheme has length two.
Consequently its codimension-two critical cycle is
\[
                              2[E_{II}].
 \tag{D.IIF3}
\]

The relative frame obstruction has the same coefficient,
without any division by two.  The normal clutch can be written
\[
 \frac1{\sqrt{4|y|^2+9|x|^4}}
   \begin{pmatrix}2y&-3\overline{x}^{\,2}\\
                   3x^2&2\overline y\end{pmatrix}\in SU(2).
 \tag{D.IIF4}
\]
Multiplication by $df$ kills its first column and sends its
second column to the positive squared denominator.
The map on the normal three-sphere has degree two: it is the
normalized map $(x,y)\mapsto(2y,3x^2)$, and a small regular
target value with second coordinate nonzero has two preimages,
both with positive complex local degree.  In the common
orientation convention of the nodal Bott matrix, this is twice
the same generator.  There are no other critical points in
this original type-$II$ chart.

For the exact fibre marking, let $B$ be the elliptic order-six
monodromy and $q_0$ its affine long-column vector.  Since
$\det(B-I)=1$, the equation $(B-I)z=q_0$ has a unique integral
solution.  In the recorded full lattice basis
$(e_1,e_2,\delta,\gamma)$ put
$\ell=(-z,0,1)$.  The full cap matrix fixes both $\ell$ and
$\delta$; its remaining row identity is part of the actual
marked cap relation.  Their principal torus at the cusp point
is exactly $E_{II}$, so the regular torus specializes with
degree one.  With its complex principal-torus orientation the
critical-cycle preimage is therefore
\[
                             2\delta\wedge\ell.
 \tag{D.IIF5}
\]
All specialization chains can be taken inside the type-$II$
cap.  This supplies the geometric multiplicity and integral
marking required by the global nullchain computation; it is
stronger than merely knowing a cohomological degree on that
curve.

\paragraph{A free good logarithmic cap: no corrected critical cycle.}
Let $\widetilde X\to\Delta_s$ be the actual smooth good-reduction
torus family used in the multiplicity-$m$ logarithmic construction,
and let its cyclic affine action have order $m=6k$, act freely
on the total space, and send $s$ to $\zeta_m s$.
These are the established hypotheses of the allowed simply
connected finite logarithmic caps; no isotriviality of the
torus family is required.  The quotient map of total spaces
$\pi:\widetilde X\to X_{\rm cap}$ is etale, while the base
coordinate is $q=s^m$.  If $D$ is the reduced central fibre,
\[
                   f^*(0)=mD,\qquad
                   df=m s^{m-1}ds
 \tag{D.IIF6}
\]
on the etale cover.  Remove the full common divisorial factor
by putting
\[
          L=f^*T\mathbb P^1\otimes\mathcal O(-(m-1)D).
 \tag{D.IIF7}
\]
The resulting differential $df_{\rm red}:TX\to L$ pulls back
to $m\,ds$.  It is surjective.  This operation is invariant
under the deck action: the local factor $s^{m-1}$ and the
transformation of $ds$ combine to the invariant original
differential.  It therefore descends as the stated bundle
map on the smooth quotient.

More generally, replacing $q$ by $u s^m$ for a holomorphic
unit $u$ gives $df_{\rm red}=mu\,ds+s\,du$, still nonzero
along $s=0$, and the nearby family is a submersion.  Thus
the exact sequence on this cap is
\[
                  0\longrightarrow V\longrightarrow TX
                      \xrightarrow{df_{\rm red}}L
                       \longrightarrow0.
 \tag{D.IIF8}
\]
Its corrected critical scheme is empty.  In particular an
uncorrected use of the vanishing $s^{m-1}$ in (D.IIF6) would
introduce a false critical contribution of codimension one.
The scalar vector is nowhere zero in the vertical bundle on
the cover and descends, since it commutes with the affine
deck action.  Hence the abstract nonzero-$\delta$ frame
argument continues across this cap.

\paragraph{Compatibility with the global relative framing calculation.}
For several multiple good fibres, use the global target line
\[
 L=f^*T\mathbb P^1\otimes
       \mathcal O\left(-\sum_i(m_i-1)D_i\right).
 \tag{D.IIF9}
\]
The kernel is unchanged by removing these common factors.
It is locally free at the original type-$II$, nodal, and
triangular charts, and a subbundle at every good logarithmic
cap.  Its abstract $\delta$ section is nowhere zero everywhere.
For the cyclic-torsion manifolds $H^2(X;\mathbb Z)=0$, so $L$
and $V/\mathbb C\delta$ are again topologically trivial.
The complete differential complex gives the literal cycle
\[
 C_{\rm crit}=\sum_{\rm nodal}E_i+
                \sum_{\rm triangular}\bigl(C_{i,1}+C_{i,2}+C_{i,3}\bigr)
                  +2\sum_{\rm original\ II}E_{II}.
 \tag{D.IIF10}
\]
There is no term at a good logarithmic cap.  Unit changes of
the line target and its divisor trivializations contribute
only $\mathbb C^*$ transition factors on a regular torus,
so they cannot change the degree-three comparison class.
After $c_2(X)=0$ has separately been established, the entire
relative-Chern/linking calculation of D.2.44 therefore applies
to (D.IIF10), with its actual marked integral preimages.
The final three-frame decision still requires a normally
framed three-cycle which is an odd generator of the
two-primary homology.  In a good logarithmic cap the ordinary
regular three-torus can become an even class, so its framing
parity alone is not sufficient.  This local lemma does not
assert that missing generator comparison.  It also makes no
assertion about the yet untreated star-fibre or nonfree-quotient
caps.
